\documentclass[12pt,a4paper,oneside]{article}
\usepackage[brazil]{babel}
\usepackage[utf8]{inputenc}
\usepackage{olymp}
\usepackage{color}
\usepackage[dvipsnames]{xcolor}
\usepackage{listings}

\setcounter{problem}{1}

\begin{document}

\begin{problem}{Inverter}{inverter}{2}
 
Dada uma lista de $N$ números inteiros, seu programa deve imprimir estes números na ordem inversa à da entrada.
%\Notes
%
%Observações, se tiver.

\InputFile

A entrada contém duas linhas: a primeira contém um valor inteiro $N$, indicando a quantidade de números na lista; a segunda contém os $N$ números da lista, que são números inteiros $V_i$ separados por espaço em branco.

Restrições: $1 \leq N \leq 1000$, $-10^6 \leq V_i \leq 10^6$.



\OutputFile

Seu programa deve gerar apenas uma linha de saída, contendo todos os $N$ números da lista, em ordem inversa.

\Notes
Escreva espaço em branco para separar os números, mas não escreva espaço em branco antes do primeiro nem depois do último número da lista. Após o último escreva apenas um caractere de quebra de linha.


\Example

\begin{example}
\exmp{
10
1 2 3 4 5 6 7 8 9 10
}{
10 9 8 7 6 5 4 3 2 1
}%
\end{example}

\begin{example}
\exmp{
7
5 7 2 14 9 -8 3
}{
3 -8 9 14 2 7 5
}%
\end{example}

\begin{example}
\exmp{
1
2
}{
2
}%
\end{example}


%Comentários sobre o exemplo, se necessário.

\end{problem}

\end{document}
